\documentclass[12pt]{article}
\usepackage[utf8]{inputenc}
\usepackage[T1]{fontenc}
\usepackage[margin=1in]{geometry}
\usepackage{setspace}
\usepackage{xcolor}
\usepackage[colorlinks=true, urlcolor=blue!60!black, linkcolor=blue!60!black]{hyperref}
\setlength{\parindent}{2em}
\setlength{\parskip}{1.25em}
\renewcommand{\baselinestretch}{1.0}
\begin{document}
\begin{center}
{\LARGE \textit{Fulton v. City of Philadelphia}, 593 U.S. 522 (2021)\par}
\vspace{0.8em}
{\large Case Brief --- Four-Page Edition\par}
\end{center}
\vspace{1em}

\section{Memory Jogger}

Unanimously reversing while declining to overrule \textit{Employment Division v. Smith}, the Court held that Philadelphia violated the Free Exercise Clause by refusing to contract with Catholic Social Services unless it certified same-sex foster couples, because a discretionary-exemption clause made the policy not generally applicable and it failed strict scrutiny (\href{https://www.courtlistener.com/opinion/4892570/fulton-v-philadelphia/}{CourtListener}).

\section{Detailed Case Facts}

Philadelphia's foster care system relies on annual contracts with private, state-licensed agencies, which certify prospective foster families through a three-to-six-month home study. Catholic Social Services (CSS), affiliated with the Archdiocese, had contracted with the City for over 50 years. Viewing certification as an endorsement of a couple's relationship, CSS would not certify unmarried couples or married same-sex couples, though it would certify single gay individuals. No same-sex couple had ever sought certification from CSS, and CSS would refer any to one of more than 20 other agencies.

In 2018, after a newspaper story, the City Council called for an investigation into ``discrimination that occurs under the guise of religious freedom.'' The DHS Commissioner told CSS that ``things have changed since 100 years ago.'' DHS then froze referrals to CSS and refused to renew its contract unless it certified same-sex couples, citing contract section 3.21, which barred rejecting foster parents based on sexual orientation ``unless an exception is granted by the Commissioner or the Commissioner's designee, in his/her sole discretion,'' and the Fair Practices Ordinance's ban on discrimination in ``public accommodations.''

\section{Procedural History}

CSS and affiliated foster parents, including Sharonell Fulton, sued under 42 U.S.C.~\S~1983. The district court denied a preliminary injunction, finding the policies neutral and generally applicable under \textit{Smith}, and the Third Circuit affirmed. The Supreme Court granted certiorari, including on whether to revisit \textit{Smith}, and reversed on June 17, 2021.

\section{Judicial Votes}

Reversed and remanded, \textbf{9--0}.
\begin{itemize}
\item \textbf{Majority:} Chief Justice Roberts, joined by Justices Breyer, Sotomayor, Kagan, Kavanaugh, and Barrett.
\item \textbf{Concurring:} Justice Barrett, joined by Justice Kavanaugh and (except the first paragraph) Justice Breyer.
\item \textbf{Concurring in the judgment:} Justice Alito, joined by Justices Thomas and Gorsuch; Justice Gorsuch, joined by Justices Thomas and Alito.
\end{itemize}

\section{Holding}

A policy is not ``generally applicable'' under \textit{Smith} if it creates a formal mechanism for individualized discretionary exemptions, whether or not any exemption has been granted. Section 3.21's ``sole discretion'' clause triggered strict scrutiny; the Fair Practices Ordinance did not apply because foster certification is not a ``public accommodation''; and the City lacked a compelling interest in denying CSS an exception. The Court declined to decide whether to overrule \textit{Smith}.

\section{Analysis of Opinions}

\textbf{Majority (Roberts).} The City burdened CSS's religious exercise by forcing it to choose between its mission and approving relationships contrary to its beliefs. Under \textit{Sherbert} and \textit{Lukumi}, an exemption mechanism invites officials to decide which reasons for noncompliance merit solicitude. Foster certification, a selective months-long evaluation, is unlike hotels, restaurants, and buses. Under \textit{O Centro}, the question was whether the City had a compelling interest in denying an exception to CSS specifically. Keeping CSS would likely help maximize foster families, liability fears were speculative, and the exemption system undermined the claim that equal treatment could brook no departures.

\textbf{Barrett (concurring).} Founding-era history is ``more silent than supportive,'' but the textual and structural case against \textit{Smith} is strong. Replacing it raises hard questions, such as treatment of institutions and indirect burdens, so she would not overrule it here.

\textbf{Alito (concurring in the judgment).} In a 77-page opinion he would overrule \textit{Smith} and restore \textit{Sherbert}. The 1791 meaning of ``prohibiting the free exercise of religion'' and early state constitutions' carveouts for ``public peace'' and ``safety'' show more than a nondiscrimination rule, and \textit{Smith} has proved unworkable. The City could nullify the ruling by deleting the clause.

\textbf{Gorsuch (concurring in the judgment).} The majority was ``dodging'' the question presented through a ``statutory shell game'' and ``a dizzying series of maneuvers''; narrow rulings force believers into endless litigation.

\section{Examples: Future Applications}

\begin{enumerate}
\item \textbf{Same-Side: ``Good Cause'' Waiver.} A homelessness agency imposes a code of conduct on faith-based shelters ``unless the Director grants a written waiver for good cause.'' The never-used waiver triggers strict scrutiny.
\item \textbf{Same-Side: Selective Licensing.} A state licenses nonprofit grief counselors after individualized vetting and allows a discretionary waiver. The process is not a public accommodation and the waiver invites discretion.
\item \textbf{Opposite-Side: Zero-Discretion Contract.} A city deletes its exception clause, leaving an absolute rule. As Alito warned, general applicability is restored under \textit{Smith}.
\item \textbf{Opposite-Side: Exception-Free Accommodations Law.} A statute bars hotels and restaurants from discriminating, with no waivers. These are squarely public accommodations, so \textit{Smith} controls.
\item \textbf{Fence-Sitter: Neutral Rule, Hostile Remarks.} A city removes all exemptions and extends its ordinance to foster agencies, but council members mock Catholic doctrine. Whether legislators' hostility defeats neutrality is unresolved under \textit{Masterpiece Cakeshop} and \textit{Lukumi}.
\end{enumerate}

\section{Critique}

\textbf{Originalist critiques.} Echoing Alito and Gorsuch, critics say \textit{Fulton} ``kicked the can down the road,'' that \textit{Smith} cannot be reconciled with original meaning, and that the majority's contorted reading of the contract and ordinance manufactured unanimity. Others praise Barrett's caution, since founding-era history is contested.

\textbf{Progressive critiques.} Ira Lupu, Robert Tuttle, and Nelson Tebbe argue that treating a never-exercised boilerplate clause as fatal stretches the individualized-exemptions doctrine and hollows out \textit{Smith} without overruling it. Critics add that the Court ignored dignitary harms to LGBTQ+ foster parents and youth when a taxpayer-funded public service turns away qualified couples.

\section{Key Quotations}

\begin{quote}
``The creation of a formal mechanism for granting exceptions renders a policy not generally applicable, regardless whether any exceptions have been given, because it `invite[s]' the government to decide which reasons for not complying with the policy are worthy of solicitude, \textit{Smith}, 494 U.S. at 884---here, at the Commissioner's `sole discretion.'\,''
\end{quote}
\noindent --- Chief Justice Roberts, Opinion of the Court.

\begin{quote}
``This decision might as well be written on the dissolving paper sold in magic shops. The City has been adamant about pressuring CSS to give in, and if the City wants to get around today's decision, it can simply eliminate the never-used exemption power.''
\end{quote}
\noindent --- Justice Alito, concurring in the judgment, 593 U.S. at 551--552.

\begin{quote}
``As a matter of text and structure, it is difficult to see why the Free Exercise Clause---lone among the First Amendment freedoms---offers nothing more than protection from discrimination.''
\end{quote}
\noindent --- Justice Barrett, concurring.

\end{document}
